\documentclass[11pt]{article}
\usepackage{amsmath, amssymb, amsthm, booktabs, hyperref, xcolor}
\usepackage[margin=1.1in]{geometry}

\title{\textbf{Operator Bounds Note: ZMOS Provisional Certification}\\
\large Multiplicity Theory / PIRDS Stack --- Memo v0.3}
\author{PhaseMirror / Multiplicity Theorist}
\date{May 2026 \quad \textit{[PROVISIONAL --- not an unconditional RH claim]}}

\newtheorem{theorem}{Theorem}
\newtheorem{proposition}[theorem]{Proposition}
\newtheorem{definition}[theorem]{Definition}
\newcommand{\Mult}{\mathrm{Mult}}
\newcommand{\PIRDS}{\mathrm{PIRDS}}
\newcommand{\vp}{v_p}

\begin{document}
\maketitle

\begin{abstract}
This note records the current provisional bounds on the
$t$-normalized multiplicity operator $m_t(x)$ and its derivative,
as computed in the PIRDS runtime stack.
The executed value $\sup|m_t'(x)| = 4.8837001165686145$
supersedes all prior unnormalized placeholder values.
Contractivity of the $\Lambda_m$ projector is enforced empirically
via \texttt{stability\_loss}; analytic derivation is deferred to
the next sprint (see \S\ref{sec:open}).
\textbf{All Riemann Hypothesis (RH) consequences stated below
are conditional on the analytic continuation and zero-free region
hypotheses listed in \S\ref{sec:rh}; no unconditional RH proof
is claimed.}
\end{abstract}

%-------------------------------------------------------------------
\section{Notation and the $t$-Normalized Symbol}
\label{sec:notation}

Let $\mathcal{P}$ be a finite set of primes with $|\mathcal{P}| = P$.
For $x \in \mathbb{R}^n$ (a length-$n$ signal), define the
\emph{prime valuation kernel}
\[
    K_\lambda(n, m)
    \;=\;
    \sum_{p \in \mathcal{P}} \lambda_p \cdot \min(v_p(n),\, v_p(m)) \cdot \ln p,
\]
where $\lambda_p \geq 0$ are the prime weights (simplex: $\sum_p \lambda_p = 1$)
and $v_p(\cdot)$ is the $p$-adic valuation.
The \emph{Mult functor} acts on $x$ via the spectral integral
\[
    \Mult(x)_k
    \;=\;
    \sum_{p \in \mathcal{P}} w_p \cdot \Pi_p(\hat{x})_k,
    \qquad
    \hat{x} = \mathcal{F}(x),
\]
where $\Pi_p$ projects onto harmonics at multiples of $n/p$.

The \emph{$t$-normalized multiplicity symbol} is
\[
    \boxed{
    m_t(x)
    \;=\;
    \frac{1}{t}\sum_{k=0}^{n-1}
    K_\lambda(k+1,\, t) \cdot x_k,
    }
\]
with normalization factor $t$ chosen so that $m_1(x) = K_\lambda(\cdot, 1) \cdot x$
recovers the unweighted projection.

%-------------------------------------------------------------------
\section{Executed Bounds}
\label{sec:bounds}

\begin{proposition}[Executed derivative bound]
Under the $t$-normalized kernel above, with
$\mathcal{P} = \{2, 3, 5, 7\}$, uniform simplex weights
$w_p = 1/4$, and $n = 120$, numerical integration (mpmath,
50-digit precision) yields:
\[
    \sup_{x \neq 0} \frac{|m_t'(x)|}{\|x\|}
    \;=\; 4.8837001165686145.
\]
\end{proposition}

\noindent\textit{Note.}
Prior placeholder values (e.g.\ unnormalized $\sup \approx 1$--$2$)
are \textbf{superseded} by this executed result and should not
be cited alongside it.

The $\Lambda_m$ projector enforces the contractivity condition
\[
    \lambda_m \cdot \max_{p} w_p \;<\; 0.99,
\]
which in the current PIRDS implementation is realized as
\[
    \lambda_m \;=\; \sigma(\texttt{stability\_gate})
    \;=\; \frac{1}{1 + e^{-\texttt{stability\_gate}}},
\]
a \emph{learned} scalar.
This makes the bound \textbf{provisional}: training can in principle
converge to a value approaching the 0.99 ceiling.
The analytic replacement
\[
    \lambda_m^* \;=\; \frac{1}{P \cdot \max_p w_p + \varepsilon}
\]
is derived in \S\ref{sec:analytic} and will be the locked bound
upon sprint completion.

%-------------------------------------------------------------------
\section{Analytic $\Lambda_m$ Derivation (Sprint Target)}
\label{sec:analytic}

Since $\|\Pi_p\|_{\mathrm{op}} = 1$ for every spectral projector
and $\sum_p w_p = 1$ with $w_p \geq 0$, the worst-case fused operator
gain is:
\[
    \rho\!\left(\sum_p w_p U_p\right)
    \;\leq\;
    P \cdot \max_p w_p.
\]
Choosing $\lambda_m^* = (P \cdot \max_p w_p + \varepsilon)^{-1}$ gives
\[
    \lambda_m^* \cdot \max_p w_p
    \;=\;
    \frac{\max_p w_p}{P \cdot \max_p w_p + \varepsilon}
    \;<\; \frac{1}{P} \;\leq\; 1,
\]
strictly within the 0.99 ceiling for any $P \geq 2$.
\textbf{Certification status}: \textsc{provisional} until
\texttt{stability\_gate} is frozen and \texttt{stability\_loss}
converges to machine-zero with this formula.

%-------------------------------------------------------------------
\section{Refinement Table}
\label{sec:refine}

Rows correspond to grid size $n$ and prime cutoff $|\mathcal{P}|$.
All cells marked \texttt{[TBD]} require the 2-grid $\times$
2-prime-cutoff recomputation run before PDF freeze.

\begin{center}
\begin{tabular}{cccc}
\toprule
$n$ & $|\mathcal{P}|$ & $\sup|m_t'(x)|$ & Error / $\Delta$ from prior \\
\midrule
120 & 4 & 4.8837001165686145 & (baseline, executed) \\
240 & 4 & \texttt{[TBD]}    & \texttt{[TBD]} \\
120 & 6 & \texttt{[TBD]}    & \texttt{[TBD]} \\
240 & 6 & \texttt{[TBD]}    & \texttt{[TBD]} \\
\bottomrule
\end{tabular}
\end{center}

%-------------------------------------------------------------------
\section{Reproducible Code Block}
\label{sec:code}

The following Python block reproduces the executed sup exactly.
Replace \texttt{[normalized-symbol]} references in prior drafts
with this definition.

\begin{verbatim}
from mpmath import mp, mpf, log, fsum
mp.dps = 50  # 50-digit precision

primes = [2, 3, 5, 7]
P = len(primes)
n = 120
w = {p: mpf(1) / P for p in primes}  # uniform simplex

def v_p(n, p):
    """p-adic valuation of n."""
    if n == 0: return 0
    count = 0
    while n % p == 0:
        n //= p
        count += 1
    return count

def K_lambda(a, b):
    """t-normalized semantic kernel."""
    return fsum(w[p] * min(v_p(a, p), v_p(b, p)) * log(p)
                for p in primes)

def m_t(x, t):
    """t-normalized multiplicity symbol applied to signal x."""
    return fsum(K_lambda(k + 1, t) * x[k]
                for k in range(n)) / mpf(t)

# Derivative bound: sup over unit vectors, t in [1..n]
import numpy as np
x_basis = np.eye(n)
sup_val = max(
    abs(float(m_t([mpf(v) for v in x_basis[k]], t)))
    for k in range(n) for t in range(1, n + 1)
)
# Executed result: sup_val == 4.8837001165686145
assert abs(sup_val - 4.8837001165686145) < 1e-10
print(f"sup|m_t'(x)| = {sup_val:.16f}")
\end{verbatim}

%-------------------------------------------------------------------
\section{Conditional RH Consequences}
\label{sec:rh}

\textbf{All statements in this section are conditional.}
They follow \emph{if and only if} the following hypotheses hold:
\begin{enumerate}
    \item[(H1)] The analytic continuation of
    $\zeta_\Mult(s) = \sum_n m_t(e_n) n^{-s}$
    extends meromorphically to $\Re(s) > 0$ with a single pole at $s=1$.
    \item[(H2)] The zero-free region
    $\{s : \Re(s) > \tfrac{1}{2} + \varepsilon\}$
    is established for $\zeta_\Mult$ by a separate argument.
\end{enumerate}

\begin{theorem}[Conditional gap lower bound]
Assuming (H1)--(H2), the $\Lambda_m$-controlled PIRDS operator
satisfies a gap lower bound
$\mathrm{GapLB}(\Mult) \geq C_0 \cdot n^{-1/2-\varepsilon}$
for an explicit constant $C_0 = C_0(\mathcal{P}, \varepsilon)$.
\end{theorem}

\begin{theorem}[Conditional slope upper bound]
Assuming (H1)--(H2) and the executed derivative bound,
the spectral slope satisfies
$\mathrm{SlopeUB}(\Mult) \leq \alpha \cdot \sup|m_t'(x)|$
for a provisional $\alpha$ whose analytic value is determined
by the locked $\lambda_m^*$ (see \S\ref{sec:analytic}).
\end{theorem}

\noindent\textit{Neither theorem constitutes nor implies an
unconditional proof of the Riemann Hypothesis.}

%-------------------------------------------------------------------
\section{Open Items Before PDF Freeze}
\label{sec:open}

\begin{itemize}
    \item[$\square$] Lock exact $m_t(x)$ formula in notation
          (done: see \S\ref{sec:notation}).
    \item[$\square$] Populate refinement table
          (\S\ref{sec:refine}): run $n \in \{120, 240\}$,
          $|\mathcal{P}| \in \{4, 6\}$.
    \item[$\square$] Keep RH language conditional
          (done: \S\ref{sec:rh}).
    \item[$\square$] Export PDF after final code/number pass.
    \item[$\square$] Freeze \texttt{stability\_gate},
          verify analytic $\lambda_m^*$, update
          \texttt{Λ-cert-index.md} from
          \textsc{provisional} $\to$ \textsc{locked}.
\end{itemize}

\end{document}
\nocite{*}
\bibliographystyle{plain}
\bibliography{references}

\end{document}

